\documentclass[../main]{subfiles}

\begin{document}
\section{Problema de Dois Corpos}\label{section:problema-de-dois-corpos}

\subsection{Preliminares}\label{subsection:fundamentacao-matematica}
\begin{defi}[Referencial Inercial]
  Um referencial inercial é um sistema de referência no qual valem as Leis de Newton. Matematicamente, ele é modelado por um espaço afim euclidiano tridimensional $E$ sobre o corpo $\mathbb{R}$, associado a um espaço vetorial de translações $V$ tal que $\dim(V)=3$, munido das seguintes estruturas:

  \begin{enumerate}
    \item Um produto interno euclidiano
      \[
        \phi: V \times V \longrightarrow \mathbb{R},
      \]
      que induz a norma
      \[
        \|v\| = \sqrt{\phi(v,v)}, \quad \forall v \in V,
      \]
      e a distância entre dois pontos $u,v \in E$ dada por
      \[
        d(u,v) = \| u - v \|.
      \]

    \item Uma orientação espacial, que juntamente com a métrica induz o produto vetorial
      \[
        \psi : V \times V \longrightarrow V.
      \]

    \item Invariância das leis da física sob transformações entre referenciais inerciais, isto é, se $S$ e $S'$ são referenciais inerciais, então as leis físicas assumem a mesma forma em ambos.

    \item Princípio da Inércia: se a força resultante sobre uma partícula é nula, então seu movimento é retilíneo e uniforme.
  \end{enumerate}
\end{defi}

\begin{defi}[Transformação de Galileu]
  Uma transformação de Galileu é uma aplicação entre dois referenciais inerciais $S$ e $S'$ no espaço-tempo clássico, definida por
  \[
    \begin{cases}
      t' = t + t_0, \\
      r' = R\,r + v\, t + r_0,
    \end{cases}
  \]
  onde $t_0 \in \mathbb{R}$ é um deslocamento temporal, $r_0 \in V$ é um vetor de translação espacial, $v \in V$ é o vetor velocidade constante relativa entre os referenciais, e $R \in SO(3)$ é uma transformação de rotação ortogonal.
\end{defi}

\begin{defi}[Grupo de Galileu]
  O grupo de Galileu é o grupo de Lie formado por todas as transformações de Galileu, munido da operação de composição de funções. Ele descreve as simetrias do espaço-tempo clássico. Seu conjunto subjacente é dado por
  \[
    G = \left\{(R,t_0,r_0,v) \;\middle|\; R \in SO(3),\; t_0 \in \mathbb{R},\; r_0 \in V,\; v \in V \right\}.
  \]
  A ação de um elemento $(R,t_0,r_0,v) \in G$ sobre um evento do espaço-tempo $(t,r)$ é dada por
  \[
    \begin{cases}
      t' = t + t_0, \\
      r' = R\,r + v\,t + r_0.
    \end{cases}
  \]
  A lei de composição em $G$ é definida por
  \[
    (R_2,t_2,r_2,v_2) \circ (R_1,t_1,r_1,v_1)
    =
    \bigl(
      R_2 R_1,\;
      t_2 + t_1,\;
      R_2 r_1 + v_2 t_1 + r_2,\;
      R_2 v_1 + v_2
    \bigr).
  \]
\end{defi}

\begin{prop}[Invariância da aceleração sob transformações de Galileu]
  Sejam \(S\) e \(S'\) dois referenciais inerciais relacionados por uma transformação de Galileu
  \[
    \begin{cases}
      t' = t + t_0,\\[4pt]
      \vec r' = R\,\vec r + \vec v\,t + \vec r_0,
    \end{cases}
  \]
  com \(R\in SO(3)\), \(\vec v,\vec r_0\in V\) e \(t_0\in\mathbb{R}\). Seja \(\vec r(t)\) a trajetória de uma partícula em \(S\). Então a aceleração \(\vec a(t)=\dfrac{d^2\vec r}{dt^2}\) transforma-se em
  \[
    \vec a'(t') = R\,\vec a(t),
  \]
  em particular, as componentes da aceleração relativas a sistemas cartesianos ortonormais são preservadas até rotação; se \(R=I\) (sem rotação) vale \(\vec a'=\vec a\).
\end{prop}

\begin{proof}
  A ação do elemento do grupo sobre a trajetória fornece, para todo \(t\),
  \[
    \vec r'(t') = R\,\vec r(t) + \vec v\,t + \vec r_0,
  \]
  onde \(t'=t+t_0\). Diferenciando em relação a \(t\) (note que \(dt'/dt=1\)) obtemos a velocidade no novo referencial:
  \[
    \vec u'(t') \equiv \frac{d\vec r'}{dt'} = \frac{d\vec r'}{dt}
    = R\,\frac{d\vec r}{dt} + \vec v = R\,\vec u(t) + \vec v.
  \]
  Derivando novamente em relação a \(t\) (ou \(t'\)) vem a aceleração:
  \[
    \vec a'(t') \equiv \frac{d^2\vec r'}{d t'^2} = \frac{d}{dt}\bigl(R\,\vec u(t) + \vec v\bigr)
    = R\,\frac{d\vec u}{dt} + \underbrace{\frac{d\vec v}{dt}}_{=0}
    = R\,\vec a(t).
  \]
  Isto prova a afirmação. Em particular, para \(R=I\) segue \(\vec a'=\vec a\).
\end{proof}

\begin{prop}[Invariância das leis de Newton (forma \(m\vec a=\vec F\)) sob o grupo de Galileu]
  Considere a segunda lei de Newton em um referencial inercial \(S\),
  \[
    m\,\vec a(t) = \vec F(t,\vec r,\dot{\vec r}),
  \]
  onde \(\vec F\) é a força resultante sobre a partícula (que pode depender de \(t,\vec r,\dot{\vec r}\) conforme o modelo). Se a força transforma-se covariantemente sob rotações e translações espaciais, i.e. sob a ação de um elemento \((R,t_0,\vec r_0,\vec v)\in G\) satisfaz
  \[
    \vec F'(t',\vec r',\dot{\vec r}') = R\,\vec F(t,\vec r,\dot{\vec r}),
  \]
  então a equação \(m\vec a=\vec F\) mantém sua forma em \(S'\):
  \[
    m\,\vec a'(t') = \vec F'(t',\vec r',\dot{\vec r}').
  \]
  Em outras palavras, a equação dinâmica é invariante sob transformações do grupo de Galileu.
\end{prop}

\begin{proof}
  Pelo Proposição anterior sabemos que \(\vec a'(t')=R\,\vec a(t)\). Assumindo que a massa \(m\) é invariante (escalar) e que a força transforma de forma covariante conforme enunciado, substituímos na esquerda e direita da segunda lei em \(S\):
  \[
    m\,\vec a'(t') = m\,(R\,\vec a(t)) = R\,(m\,\vec a(t)) = R\,\vec F(t,\vec r,\dot{\vec r}) = \vec F'(t',\vec r',\dot{\vec r}'),
  \]
  o que demonstra que a forma da lei é preservada nas transformações de Galileu.
\end{proof}

\paragraph{Comentários adicionais.}
\begin{itemize}
  \item A hipótese de covariância da força é natural para forças de origem geométrica (por exemplo forças derivadas de potenciais escalares \(V(\vec r)\) que satisfazem \(V'( \vec r') = V(\vec r)\) quando \(R\) é uma rotação e as translações são homogêneas). Para forças dependentes da velocidade é necessária a transformação adequada de \(\dot{\vec r}\) (por ex. \(\dot{\vec r}'=R\dot{\vec r}+\vec v\)).
  \item A invariância da primeira lei (princípio da inércia) decorre imediatamente: se em \(S\) a força resultante é nula então \( \vec a=0\); aplicando a proposição obtemos \(\vec a'=R\vec a=0\), logo em \(S'\) a aceleração também é nula.
  \item A terceira lei (ação e reação) exige hipóteses adicionais sobre a forma das forças internas (ex.: forças centrais iguais e opostas), mas rotações e translações preservam a relação ação–reação quando as forças transformam covariantemente por rotação.
\end{itemize}

\end{document}
